%! Orthogonal Matrices and Gram--Schmidt — Unit 4, Lesson 4.4 — Guided Notes (Answer Key)
\documentclass[10pt]{article}
\usepackage{linalg-article}
\usepackage{linalg-key}   % pulls in -boxes; do NOT also load -boxes
\newcommand{\TallMath}[1]{$\displaystyle #1\rule[-1.4em]{0pt}{3.2em}$}
\newcommand{\vv}{\mathbf{v}}
\newcommand{\ww}{\mathbf{w}}
\newcommand{\xx}{\mathbf{x}}
\newcommand{\bb}{\mathbf{b}}
\newcommand{\pp}{\mathbf{p}}
\newcommand{\avec}{\mathbf{a}}
\newcommand{\qq}{\mathbf{q}}
\newcommand{\zero}{\mathbf{0}}
\newcommand{\T}{\mathsf{T}}

\begin{document}

\pageheader{Unit 4, Lesson 4.4}{Guided Notes --- Answer Key}
\namedateperiod

\vspace{0.10in}

\begin{objectivebox}
\textbf{By the end of this lesson, I will be able to\dots}
\begin{itemize}\setlength{\itemsep}{2pt}
  \item test whether vectors are \emph{orthonormal} (perpendicular and unit length) and form the matrix $Q$ with $Q^{\T}Q=I$;
  \item find coordinates and projections with no solving --- each weight is a dot product, so $\hat{\xx}=Q^{\T}\bb$;
  \item run \emph{Gram--Schmidt} to turn independent vectors into an orthonormal set, and say why it works.
\end{itemize}
\end{objectivebox}

\vspace{0.04in}

\begin{vocabbox}
\small
Fill in each term as we build it together.
\vspace{2pt}
\vocabans{Orthonormal vectors}{vectors that are mutually perpendicular ($q_i\cdot q_j=0$ for $i\ne j$) and each of unit length ($q_i\cdot q_i=1$)}
\vocabans{Orthogonal matrix $Q$ ($Q^{\T}Q=I$)}{a matrix whose columns are orthonormal; when it is square, $Q^{\T}=Q^{-1}$}
\vocabans{Gram--Schmidt process}{a procedure turning independent vectors into orthonormal ones by subtracting off overlaps, then normalizing}
\end{vocabbox}

\begin{hookbox}
Compare two coordinate grids: ordinary graph paper (square, perpendicular lines, unit spacing) and a
skewed, stretched grid.
\begin{itemize}\setlength{\itemsep}{1pt}
  \item On \emph{graph paper}, how do you find a point's coordinates? \ansline{Read straight across and up --- the axes are perpendicular and unit-spaced, so each coordinate is immediate, no computation.}
  \item On the \emph{skewed} grid, why is it more work? \ansline{The axes are not perpendicular (or not unit length), so you must \emph{solve} for the weights that combine them --- exactly the trouble orthonormal axes remove.}
\end{itemize}
\end{hookbox}

\begin{notesbox}{1. Orthonormal Vectors}
``Ortho'' means \emph{perpendicular}; ``normal'' means \emph{length one}. A set $q_1,q_2,\dots$ is
\textbf{orthonormal} when
\[
  q_i\cdot q_j={\color{keyred}\mathbf{0}}\ \ (\text{for } i\ne j)\qquad\text{and}\qquad q_i\cdot q_i={\color{keyred}\mathbf{1}}.
\]
Take the perpendicular pair from Lesson~4.0: $(3,4)$ and $(4,-3)$. They \emph{are} perpendicular, but each
has length $5$, so they are not yet orthonormal. \textbf{Normalize} --- divide each by its length:
\[
  q_1=\tfrac15\begin{bsmallmatrix} 3\\ 4 \end{bsmallmatrix},\qquad
  q_2=\tfrac15\begin{bsmallmatrix} 4\\ -3 \end{bsmallmatrix}.
\]
Now $q_1\cdot q_1=\tfrac{1}{25}(9{+}16)=1$, $q_2\cdot q_2=1$, and $q_1\cdot q_2=\tfrac{1}{25}(12{-}12)=0$.
Perpendicular \emph{and} unit length --- perfect ``graph-paper'' axes.

\begin{center}
\begin{tikzpicture}[scale=1.0, >=stealth]
  \draw[->, charcoal] (-0.4,0)--(2.6,0) node[right, font=\scriptsize]{};
  \draw[->, charcoal] (0,-1.9)--(0,2.4) node[above, font=\scriptsize]{};
  % q1 direction (0.6,0.8) scaled 2.4 -> (1.44,1.92); q2 dir (0.8,-0.6) scaled 2.4 -> (1.92,-1.44)
  \draw[burgundy, very thick, ->] (0,0)--(1.44,1.92) node[above right, font=\scriptsize]{$q_1=\tfrac15(3,4)$};
  \draw[royalblue, very thick, ->] (0,0)--(1.92,-1.44) node[below right, font=\scriptsize]{$q_2=\tfrac15(4,-3)$};
  % right-angle square: A=0.35*q1dir=(0.21,0.28), B=0.35*q2dir=(0.28,-0.21), corner=A+B=(0.49,0.07)
  \draw[charcoal, thin] (0.21,0.28)--(0.49,0.07)--(0.28,-0.21);
  \node[charcoal, font=\scriptsize] at (2.15,0.55) {both length $1$};
\end{tikzpicture}
\end{center}
\end{notesbox}

\begin{notesbox}{2. The Matrix $Q$, and Coordinates for Free}
Stack the orthonormal vectors as the columns of a matrix $Q=\begin{bsmallmatrix} q_1 & q_2 \end{bsmallmatrix}
=\tfrac15\begin{bsmallmatrix} 3 & 4\\ 4 & -3 \end{bsmallmatrix}$. Row $i$ of $Q^{\T}$ dotted with column $j$
of $Q$ is exactly $q_i\cdot q_j$, so
\[
  Q^{\T}Q=\begin{bsmallmatrix} q_1\cdot q_1 & q_1\cdot q_2\\ q_2\cdot q_1 & q_2\cdot q_2 \end{bsmallmatrix}
  =\begin{bsmallmatrix} {\color{keyred}\mathbf{1}} & {\color{keyred}\mathbf{0}}\\ {\color{keyred}\mathbf{0}} & {\color{keyred}\mathbf{1}} \end{bsmallmatrix}=I.
\]
When $Q$ is square, $Q^{\T}Q=I$ says $Q^{\T}=Q^{-1}$: \emph{the inverse is free.} And because the axes are
orthonormal, the coordinates of any $\bb$ are just dot products:
\[
  \bb=(q_1\cdot\bb)\,q_1+(q_2\cdot\bb)\,q_2.
\]
\textbf{Try $\bb=(5,0)$:} $q_1\cdot\bb=\tfrac15(15)={\color{keyred}\mathbf{3}}$ and $q_2\cdot\bb=\tfrac15(20)={\color{keyred}\mathbf{4}}$,
so $\bb=3\,q_1+4\,q_2$. No system to solve --- read the coordinates straight off the dot products.
\vspace{-2pt}
\ansline{Check: $3\,q_1+4\,q_2=\tfrac15(9,12)+\tfrac15(16,-12)=\tfrac15(25,0)=(5,0)=\bb$. \checkmark}
\end{notesbox}

\begin{notesbox}{3. Least Squares Becomes Trivial}
Last lesson we projected by solving the normal equations $A^{\T}A\hat{\xx}=A^{\T}\bb$. If the columns of $A$
are orthonormal --- call the matrix $Q$ --- then $A^{\T}A=Q^{\T}Q=I$, and the normal equations collapse:
\[
  \boxed{\,\hat{\xx}=Q^{\T}\bb\,}\qquad\text{and}\qquad \pp=Q\hat{\xx}=QQ^{\T}\bb.
\]
Each entry of $\hat{\xx}$ is one dot product $q_i\cdot\bb$ --- \emph{there is nothing to invert.} This is why
orthonormal columns are so valuable: projection, least squares, and finding coordinates all become a handful
of dot products.
\end{notesbox}

\begin{notesbox}{4. Gram--Schmidt --- Building Orthonormal Columns}
What if our vectors are independent but \emph{not} perpendicular? \textbf{Gram--Schmidt} fixes that. Start
with independent $\avec_1=(3,4)$ and $\avec_2=(1,0)$.
\begin{itemize}\setlength{\itemsep}{3pt}
  \item \textbf{Step 1 --- normalize the first:} $q_1=\dfrac{\avec_1}{\|\avec_1\|}=\tfrac15(3,4)$.
  \item \textbf{Step 2 --- subtract the overlap of $\avec_2$ with $q_1$:}
  \[
    B=\avec_2-(q_1\cdot\avec_2)\,q_1
     =\begin{bsmallmatrix} 1\\ 0 \end{bsmallmatrix}-\tfrac35\!\cdot\!\tfrac15\begin{bsmallmatrix} 3\\ 4 \end{bsmallmatrix}
     =\begin{bsmallmatrix} 16/25\\ -12/25 \end{bsmallmatrix},\qquad \|B\|=\tfrac{4}{5}.
  \]
  \item \textbf{Step 3 --- normalize the remainder:} $q_2=\dfrac{B}{\|B\|}=\tfrac15(4,-3)$.
\end{itemize}
\textbf{Why does the subtraction work?} Check $q_1\cdot B=q_1\cdot\avec_2-(q_1\cdot\avec_2)(q_1\cdot q_1)
=\tfrac35-\tfrac35\cdot 1={\color{keyred}\mathbf{0}}$. We removed \emph{exactly} the part of $\avec_2$ pointing along
$q_1$ (a projection!), so what is left is perpendicular. The result: $Q=\tfrac15\begin{bsmallmatrix} 3&4\\4&-3 \end{bsmallmatrix}$
--- orthonormal columns, ready to use.
\end{notesbox}

\begin{practicebox}
Show your work.
\begin{enumerate}[leftmargin=1.6em, itemsep=6pt]
  \item Is the pair $q_1=\tfrac13(2,2,1)$, $q_2=\tfrac13(2,-1,-2)$ orthonormal? Compute the three dot
        products $q_1\cdot q_1$, $q_2\cdot q_2$, $q_1\cdot q_2$ and decide.
        \ansline{$q_1\cdot q_1=\tfrac19(4{+}4{+}1)=1$; $q_2\cdot q_2=\tfrac19(4{+}1{+}4)=1$; $q_1\cdot q_2=\tfrac19(4{-}2{-}2)=0$. Both unit length and perpendicular, so \emph{yes}, the pair is orthonormal.}
  \item With $Q=\tfrac15\begin{bsmallmatrix} 3&4\\4&-3 \end{bsmallmatrix}$, write $\bb=(10,5)$ in the
        orthonormal basis using dot products (find the two weights).
        \ansline{$q_1\cdot\bb=\tfrac15(30{+}20)=10$ and $q_2\cdot\bb=\tfrac15(40{-}15)=5$, so $\bb=10\,q_1+5\,q_2$. (Check: $10\,q_1+5\,q_2=(6,8)+(4,-3)=(10,5)$. \checkmark)}
\end{enumerate}
\end{practicebox}

\begin{teachernote}
The one idea: \textbf{orthonormal columns make $A^{\T}A=I$}, so every hard step from Units~2--4 (inverting,
solving the normal equations) becomes a few dot products. Build the vocabulary carefully --- students
routinely check \emph{perpendicular} and forget \emph{unit length}; stress both halves of ``orthonormal.''
The normalize step is Lesson~1.2 ($\vv/\|\vv\|$); the perpendicular test is Lesson~1.2/4.0; the projection
inside Gram--Schmidt is Lesson~4.2. The subtraction $B=\avec_2-(q_1\cdot\avec_2)q_1$ removes exactly the
component of $\avec_2$ along $q_1$ --- that is why $q_1\cdot B=0$, and it is worth showing algebraically once
(as in \S4). Common slips: dividing by the wrong length; normalizing before subtracting; treating a
merely-perpendicular pair as orthonormal. This lesson closes Unit~4 --- next is the unit test, then Unit~5
(determinants).
\end{teachernote}

\end{document}
